\chapter{Demand}
\label{chap:demand}

\section{Learning objectives}

After completing this chapter, you should be able to:
\begin{itemize}
  \item derive the idea of individual demand from consumer choice;
  \item distinguish movements along demand from shifts of demand;
  \item explain normal, inferior, and Giffen goods;
  \item distinguish income and substitution effects; and
  \item construct market demand from individual demands.
\end{itemize}

\section{From choice to demand}

An \keyterm{individual demand function} describes the quantity of a good a consumer chooses at each possible price, holding other relevant factors constant. It comes from solving the consumer's optimisation problem repeatedly as the price changes. Demand is therefore a relationship between price and the utility-maximising quantity, not simply a desire for a product.

For a consumer choosing goods $x$ and $y$, demand for $x$ can be written as
\[
x=x(p_x,p_y,I;\text{preferences}).
\]
The semicolon reminds us that preferences are treated as given in the analysis. A change in $p_x$ generates a movement along the demand curve; a change in income, the price of $y$, or preferences can shift the entire demand curve.

\section{The law of demand}

The \keyterm{law of demand} states that, other things equal, a decrease in the price of a good increases the quantity demanded. Graphically, an ordinary demand curve slopes downward. The law is not merely an empirical regularity: it follows from two forces created by a price change.

First, when $p_x$ falls, $x$ becomes relatively cheaper than $y$. The consumer substitutes toward $x$. Second, the lower price increases purchasing power: the consumer can now afford more bundles than before. For a normal good, this income effect also increases demand for $x$.

\section{Movements and shifts}

A \keyterm{movement along the demand curve} occurs when the price of the good itself changes. If the price of coffee falls, consumers move to a point with a higher quantity demanded on the same demand curve.

A \keyterm{shift in demand} occurs when another determinant changes. Important determinants include income, prices of related goods, tastes, expectations, the number of consumers, and institutions. If consumers' income rises and coffee is a normal good, the demand curve for coffee shifts rightward: at every given price, more coffee is demanded.

Two goods are \keyterm{substitutes} when an increase in the price of one raises demand for the other, such as tea and coffee for some consumers. They are \keyterm{complements} when an increase in the price of one reduces demand for the other, such as printers and printer cartridges.

\section{Income effects and Engel curves}

The relationship between income and the quantity demanded of a good is summarised by an \keyterm{Engel curve}. A good is \keyterm{normal} if demand rises when income rises, other things equal. A good is \keyterm{inferior} if demand falls when income rises.

Inferior does not mean low quality. It means that, as consumers become richer, they choose less of the good because they switch to alternatives they value more. For example, a student may reduce consumption of inexpensive ready-made meals as income rises and buy more fresh food instead.

The classification can depend on the income range. A good may be normal at low income levels and inferior at higher ones. Demand analysis is always conditional on the consumer, the alternatives available, and the relevant range of income.

\section{Substitution and income effects}

When $p_x$ changes, the consumer's choice changes for two distinct reasons. The \keyterm{substitution effect} is the change in consumption caused by the change in the relative price, holding utility or purchasing power constant according to the method used. When $x$ becomes cheaper relative to $y$, this effect increases demand for $x$.

The \keyterm{income effect} is the change caused by the resulting change in real purchasing power. A fall in $p_x$ makes the consumer effectively richer. For a normal good, this raises consumption of $x$; for an inferior good, it lowers consumption of $x$.

For most goods, the substitution effect dominates any opposing income effect, so demand slopes downward. A \keyterm{Giffen good} is a rare case in which a sufficiently strong negative income effect outweighs the substitution effect, causing quantity demanded to rise when price rises. The concept is theoretically important, but it should not be confused with a conspicuous or fashionable good.

\section{Market demand}

\keyterm{Market demand} is the total quantity demanded by all consumers at each price. It is obtained by horizontally adding individual demand curves: for a given price, add the quantities each consumer chooses.

If Ana demands $q_A=20-p$ and Bruno demands $q_B=10-p/2$ at prices for which both quantities are positive, market demand is
\[
Q=q_A+q_B=30-\frac{3}{2}p.
\]
Market demand depends not only on individual preferences and incomes but also on the number and composition of consumers in the market.

\section{Demand and empirical evidence}

Economic theory predicts how demand responds when other factors are held fixed. In practice, many factors change simultaneously. Observing that high-price products sell large quantities does not by itself disprove the law of demand; perhaps those products also have stronger demand because of quality, popularity, or income differences. Identifying causal effects requires careful evidence and institutional knowledge.

\section*{Chapter Review}

\subsection*{Key concepts}
\begin{itemize}
  \item \textbf{Individual demand}: a consumer's utility-maximising quantity at each price.
  \item \textbf{Movement along demand}: caused by a change in the good's own price.
  \item \textbf{Demand shift}: caused by a change in another determinant.
  \item \textbf{Normal good}: demand rises with income.
  \item \textbf{Inferior good}: demand falls with income.
  \item \textbf{Substitution effect}: response to a change in relative prices.
  \item \textbf{Income effect}: response to a change in purchasing power.
  \item \textbf{Market demand}: horizontal sum of individual demands.
\end{itemize}

\section*{Exercises}

\subsection*{Conceptual questions}
\begin{enumerate}[label=\arabic*.]
  \item Define demand. Why is demand not the same as a consumer's desire for a good?
  \item Distinguish a movement along a demand curve from a shift of demand.
  \item Define normal and inferior goods. Why is ``inferior'' not a moral judgement about quality?
  \item Explain the substitution and income effects of a fall in the price of a normal good.
  \item Why are Giffen goods expected to be unusual?
\end{enumerate}

\subsection*{Graphical exercises}
\begin{enumerate}[resume,label=\arabic*.]
  \item Draw a downward-sloping demand curve. Show the change in quantity demanded following a fall in price.
  \item Draw a demand curve for a normal good and show the effect of an increase in income.
  \item Draw Engel curves for a normal good, an inferior good, and a good that is normal at low income but inferior at high income.
\end{enumerate}

\subsection*{Analytical exercises}
\begin{enumerate}[resume,label=\arabic*.]
  \item Ana's demand is $q_A=16-2p$ and Bruno's demand is $q_B=10-p$, provided quantities are non-negative. Derive market demand and state the relevant price range.
  \item A consumer's demand for bus journeys falls when income increases, holding all prices constant. Classify bus journeys for this consumer and give a plausible explanation.
  \item The price of tea rises and demand for coffee rises. What does this suggest about the relationship between tea and coffee? What additional evidence would you want before reaching a conclusion?
\end{enumerate}

\subsection*{Exam-style multiple-choice questions}
\begin{enumerate}[resume,label=\arabic*.]
  \item A movement along a demand curve is caused by: (a) a change in income; (b) a change in the good's own price; (c) a change in tastes; or (d) a change in the price of a substitute.
  \item An inferior good is one whose demand: (a) rises when income rises; (b) falls when income rises; (c) rises when its own price rises; or (d) is always zero.
  \item Market demand is found by: (a) vertically adding individual demand curves; (b) averaging individual prices; (c) horizontally adding individual quantities at each price; or (d) choosing the demand of the richest consumer.
\end{enumerate}

\subsection*{PPE discussion question}

\noindent Public statistics often report correlations between income, prices, and consumption. Discuss why a causal claim about demand requires more than observing that two variables move together. Give an example relevant to public policy.
